\documentclass[11pt]{article}
\usepackage[utf8]{inputenc}
\usepackage[T1]{fontenc}
\usepackage{lmodern,microtype,booktabs,array,longtable,setspace,url,graphicx}
\usepackage[margin=0.8in]{geometry}
\usepackage[hidelinks]{hyperref}
\doublespacing
\title{Supplementary Material\\Governing Interruptions After Deterministic Deterioration Scoring}
\author{Satya Venkata Ranga Janaki Sriram Mentey}
\date{}

\begin{document}
\maketitle

\section*{Supplementary methods}

\subsection*{Frozen replay and version boundary}
The manuscript reports the Challenge 2019 study frozen at Git commit
\url{https://github.com/sreeram843/curie-prediction-pipeline/commit/1d8a1117151619e2d34822e365d17a00e30e5233}.
The setB replay was repeated
from that revision with the frozen winner and reproduced the published primary
counts. Later scorer revisions are intentionally outside this study. The
aggregate v3 holdout sidecar adds denominators, patient-time, stay-level PPV,
lead-time quartiles, alert-count distributions, and normalized Challenge utility
without replacing v1 or v2 frozen artifacts. Patient-level rows are not stored
in the repository.

\subsection*{Partial-SOFA mapping}
At each hourly row, the replay retained the most recently observed value for
each available variable within that stay; it did not use future rows. No maximum
carry-forward age was imposed in the frozen revision. O2Sat was divided by
FiO$_2$, with SaO2 used when O2Sat was absent. The resulting S/F-like ratio was
compared directly with the numeric P/F bands; no validated S/F conversion was
applied, so this is a nonstandard respiratory proxy. Because
mechanical ventilation was unavailable, respiratory scores of 3 or 4 were not
assigned. The cardiovascular component used MAP only. Renal scoring used
creatinine only. GCS and urine-output windows were unavailable. At least two
scoreable components were required, and threshold-only emission occurred at
partial SOFA $\geq2$.

\begin{table}[htbp]
\centering\small
\caption{Supplementary Table S1. Component mapping in the frozen Challenge
replay. Values at a boundary enter the less severe band unless stated.}
\begin{tabular}{@{}lll@{}}
\toprule
Component & Available input & Points \\
\midrule
Respiration & O2Sat/FiO$_2$ (SaO2 fallback) & $<300$: 2; $300$--$<400$: 1; $\geq400$: 0 \\
Coagulation & Platelets ($10^9$/L) & $<20$: 4; $<50$: 3; $<100$: 2; $<150$: 1; else 0 \\
Liver & Bilirubin (mg/dL) & $\geq12$: 4; $\geq6$: 3; $\geq2$: 2; $\geq1.2$: 1; else 0 \\
Cardiovascular & MAP (mmHg) & $<70$: 1; otherwise 0 \\
Renal & Creatinine (mg/dL) & $\geq5$: 4; $\geq3.5$: 3; $\geq2$: 2; $\geq1.2$: 1; else 0 \\
CNS & GCS & Unavailable; missing \\
\bottomrule
\end{tabular}
\end{table}

\subsection*{Selection procedure}
The original 23-candidate setA sweep used the legacy grace-6 definition: first
emission at or before \texttt{label\_start}+6 hours. It was not the later
bounded $-12/+6$-hour primary-window definition. Passing required governed
sensitivity at least naive sensitivity minus 0.10 or at least 0.70, plus an
interruptive-to-naive emission ratio no greater than 0.25. Ranking prioritized
passing both objectives, lower interruptive ratio, lower emission burden per
detected positive stay, and higher governed sensitivity. The winner had 88.4\%
governed sensitivity under the legacy grace-6 selection estimand. Recalculation
under the later primary window yielded 87.6\%; these values answer different
estimands and are not interchangeable.

\begin{table}[htbp]
\centering\scriptsize
\setlength{\tabcolsep}{2.5pt}
\caption{Supplementary Table S2. Complete 23-candidate setA selection grid under
the legacy grace-6 estimand. Persist and Refrac are minutes; Base and Gate are
baseline and page-gate states.}
\resizebox{\textwidth}{!}{\begin{tabular}{@{}lrrrrrrrr@{}}
\toprule
Candidate & Persist & Refrac & Base & Gate & Gov. sens. (\%) & Int. sens. (\%) & Int. ratio & Selected \\
\midrule
grid\_p0\_r90\_b0 & 0 & 90 & off & on & 88.4 & 37.7 & 0.123 & yes \\
grid\_p0\_r120\_b0 & 0 & 120 & off & on & 88.4 & 37.7 & 0.123 & no \\
dual & 0 & 90 & off & on & 88.4 & 37.7 & 0.123 & no \\
grid\_p15\_r90\_b0 & 15 & 90 & off & on & 86.1 & 38.6 & 0.123 & no \\
grid\_p15\_r120\_b0 & 15 & 120 & off & on & 86.1 & 38.6 & 0.123 & no \\
grid\_p30\_r90\_b0 & 30 & 90 & off & on & 86.1 & 38.6 & 0.123 & no \\
grid\_p30\_r120\_b0 & 30 & 120 & off & on & 86.1 & 38.6 & 0.123 & no \\
grid\_p0\_r60\_b0 & 0 & 60 & off & on & 88.4 & 40.3 & 0.245 & no \\
grid\_p30\_r60\_b0 & 30 & 60 & off & on & 86.1 & 40.3 & 0.245 & no \\
grid\_p15\_r60\_b0 & 15 & 60 & off & on & 86.1 & 40.3 & 0.245 & no \\
sensitive & 0 & 120 & off & off & 88.4 & 44.1 & 0.161 & no \\
accuracy & 0 & 90 & off & off & 88.4 & 44.1 & 0.161 & no \\
strict & 30 & 120 & on & off & 16.9 & 16.9 & 0.024 & no \\
grid\_p15\_r120\_b1 & 15 & 120 & on & on & 26.4 & 26.2 & 0.042 & no \\
grid\_p15\_r90\_b1 & 15 & 90 & on & on & 26.4 & 26.2 & 0.042 & no \\
grid\_p30\_r90\_b1 & 30 & 90 & on & on & 26.4 & 26.2 & 0.042 & no \\
grid\_p30\_r120\_b1 & 30 & 120 & on & on & 26.4 & 26.2 & 0.042 & no \\
grid\_p0\_r120\_b1 & 0 & 120 & on & on & 29.9 & 29.7 & 0.049 & no \\
grid\_p0\_r90\_b1 & 0 & 90 & on & on & 29.9 & 29.7 & 0.049 & no \\
balanced & 15 & 60 & on & on & 26.4 & 26.2 & 0.075 & no \\
grid\_p15\_r60\_b1 & 15 & 60 & on & on & 26.4 & 26.2 & 0.075 & no \\
grid\_p30\_r60\_b1 & 30 & 60 & on & on & 26.4 & 26.2 & 0.075 & no \\
grid\_p0\_r60\_b1 & 0 & 60 & on & on & 29.9 & 29.7 & 0.089 & no \\
\bottomrule
\end{tabular}
}
\end{table}

\subsection*{Timing sensitivity}
The primary setB estimand required any alert in the bounded interval
$[\texttt{label\_start}-12\,\mathrm{h},\texttt{label\_start}+6\,\mathrm{h}]$.
The following secondary definitions were computed without changing the frozen
governance configuration.

\begin{table}[htbp]
\centering\small
\caption{Supplementary Table S3. Secondary setB timing definitions.}
\begin{tabular}{@{}lcc@{}}
\toprule
Definition & Naive sensitivity (\%) & Governed sensitivity (\%) \\
\midrule
first alert at or before label\_start & 56.3 & 56.3 \\
first alert at or before label\_start + 6 hours & 81.1 & 81.1 \\
first alert at or before label\_start + 12 hours & 84.6 & 84.6 \\
first alert before label\_start & 54.8 & 54.8 \\
any alert in [label\_start - 12 hours, label\_start + 12 hours] & 83.3 & 83.3 \\
\bottomrule
\end{tabular}

\end{table}

\subsection*{Statistical definitions}
Sensitivity was the proportion of label-positive stays with at least one
qualifying emission in the selected window. Emissions per detected positive
stay divided all emissions in a lane by positive stays detected by that lane;
it is not a number of clinically adjudicated pages. Stay-level PPV divided
positive stays detected in-window by all stays with any emission. For negative
stays, an emission at any recorded hour counted in the denominator. Lead time
was \texttt{label\_start} minus the first governed emission within the primary
window. Confidence intervals used 1,000 stay-level percentile-bootstrap
replicates with seed 42.

Normalized Challenge utility used emit-hour predictions rather than latching a
patient positive after the first emission. It is reported as an additional
within-resource metric, not as a hidden-test leaderboard result and not as a
clinical utility estimate.

\section*{Supplementary reproducibility assets}
The frozen operating point, resolved study bundle, content hash, timing policy,
complete aggregate selection grid, holdout sidecars, comparator results,
ablations, miss attribution, bootstrap intervals, and machine-readable tables
are included in the repository. The public dataset is available at
\url{https://doi.org/10.13026/v64v-d857}. Rebuild commands and pinned hashes are
listed in \texttt{paper/REPRODUCIBILITY.md} and the manuscript manifest.

\end{document}
